%% MethodsX methods article -- Elsevier elsarticle template
%% Structure follows the official "MethodsX methods article template
%% Version 7 (September 2024)". Section headings match the template and must
%% not be changed. All sections are mandatory unless marked OPTIONAL.
%%
%% Compile with pdfLaTeX (run twice to resolve \cite references).
%% Overleaf: New Project -> Upload Project -> Compiler = pdfLaTeX.
%%
%% LAYOUT: 'final,3p' = journal single-column look (for reading / advisor
%%   review). When you SUBMIT, change 'final' to 'review' for the
%%   double-spaced, line-numbered manuscript format Elsevier expects.
%%
%% TODO before submitting:
%%   - fill every [bracketed] placeholder (affiliations, ethics,
%%     funding, CRediT roles);
%%   - prepare the Graphical abstract as a SEPARATE image file;
%%   - delete this comment header.

\documentclass[final,3p,times]{elsarticle}

\usepackage{lineno,hyperref}
\usepackage{tabularx}
\usepackage{booktabs}
\usepackage{amssymb}
\usepackage[utf8]{inputenc}
\usepackage[T1]{fontenc}
\modulolinenumbers[5]

\journal{MethodsX}

\begin{document}

%% =====================================================================
%% ARTICLE INFORMATION
%% Title (max. 20 words), authors (mark corresponding with *), full
%% affiliations, corresponding e-mail (+ X handle), keywords (>= 3).
%% =====================================================================
\begin{frontmatter}

\title{A transfer-learning and automatic false-positive cleanup pipeline
for detecting primate vocalizations in long tropical-forest recordings}

\author[a]{Moshi Fu\corref{cor1}}
\ead{mfu39@wisc.edu}
\cortext[cor1]{Corresponding author.}

\author[a]{[Co-author Name]}

\affiliation[a]{organization={[Department], University of Wisconsin--Madison},
            addressline={[Street address]},
            city={Madison},
            postcode={[ZIP]},
            state={WI},
            country={USA}}

\affiliation[b]{organization={[Second institution, full address]},
            country={[Country]}}

%% ABSTRACT -- max. 200 words, ending with 1-3 bullet points that
%% outline the method (no technical steps).
\begin{abstract}
Passive acoustic monitoring (PAM) produces thousands of hours of
tropical-forest audio that are infeasible to annotate by hand. We describe an
end-to-end pipeline that detects primate vocalizations --- the putty-nosed
monkey (\textit{Cercopithecus nictitans}) and \textit{Colobus guereza} ---
in long field recordings from Makokou, Gabon. A three-class model (Cernic,
Colobus, Background) is built by pooling all putty-nosed call types at the
training level, eliminating the need for softmax summation at detection time.
Two-second audio windows --- with short clips embedded in real background noise
and longer clips trimmed to their highest-energy window --- are converted to
mel-spectrogram images and classified by a temporal-frequency CRNN head
attached to a VGG19 backbone: the feature map is split into four frequency
bands, processed by per-band and cross-band convolutions, and fed through a
bidirectional LSTM to capture spectro-temporal structure. Two-stage training
(frozen base then partial fine-tuning) with human-verified, label-audited
training data achieves 99.1\% validation accuracy with zero cross-species
confusion. An automatic three-filter false-positive cleanup flags suspicious
detections, and confirmed false positives are recycled as hard negatives for
iterative retraining.

\begin{itemize}
\item A temporal-frequency CRNN head on a VGG19 backbone for three-class
primate call detection with two-stage transfer learning.
\item Human-verified training data with hard-negative recovery from
false-positive pools, eliminating label noise and silence shortcuts.
\item An automatic three-filter cleanup pipeline that feeds confirmed false
positives back as hard negatives for iterative retraining.
\end{itemize}
\end{abstract}

%% KEYWORDS -- minimum three, separated by \sep.
\begin{keyword}
Bioacoustics \sep Passive acoustic monitoring \sep Transfer learning \sep
Convolutional recurrent neural network \sep Mel-spectrogram \sep
False-positive filtering \sep Out-of-distribution detection \sep
\textit{Cercopithecus nictitans} \sep \textit{Colobus guereza}
\end{keyword}

\end{frontmatter}

%% =====================================================================
\section*{Related research article}
None.

%% =====================================================================
\section*{Graphical abstract}
\textit{[Submitted as a separate image file (min.\ 531 $\times$ 1328 px,
h $\times$ w; readable at 5 $\times$ 13 cm; TIFF/EPS/PDF/MS Office). Suggested
content: the pipeline --- reference clips $\rightarrow$
mel-spectrogram + temporal-frequency CRNN classifier $\rightarrow$
sliding-window detection $\rightarrow$ three-filter cleanup $\rightarrow$
hard-negative retraining loop.]}

%% =====================================================================
\section*{Specifications table}

\noindent
\renewcommand{\arraystretch}{1.3}
\begin{tabularx}{\linewidth}{@{}lX@{}}
\toprule
Subject area & Agricultural and Biological Sciences \\
More specific subject area & Bioacoustics; passive acoustic monitoring;
conservation biology \\
Name of your method & VGG19 temporal-frequency CRNN with two-stage
fine-tuning and three-filter automatic false-positive cleanup for primate
vocalizations \\
Name and reference of original method & Transfer learning with the VGG19
architecture \cite{simonyan2015}; out-of-distribution detection via the
Mahalanobis distance \cite{lee2018}; audio event tagging with YAMNet trained
on AudioSet \cite{gemmeke2017,hershey2017} \\
Resource availability & Code:
\url{https://github.com/mo119m/primates-sound-detection} --- Python 3,
TensorFlow 2, librosa, scikit-learn, OpenCV, soundfile; the YAMNet filter
additionally needs tensorflow-hub and resampy. [Add a data DOI/repository
here if you deposit the reference clips.] \\
\bottomrule
\end{tabularx}

%% =====================================================================
%% BACKGROUND -- motivation and context. Max. 500 words.
%% =====================================================================
\section*{Background}

Passive acoustic monitoring is now a standard tool for surveying vocal animals
across large areas and long time spans, but it generates volumes of audio that
cannot realistically be reviewed by ear \cite{stowell2022}. Deep convolutional
networks applied to spectrogram ``images'' have become the dominant approach
for automatically detecting species-specific calls, and transfer learning from
large image or audio corpora lets such detectors be trained from comparatively
few annotated examples \cite{dufourq2022,kahl2021}. These approaches work well
on clean, isolated reference clips.

The difficulty in practice is precision rather than recall. A model trained on
curated clips of a target call will reliably fire on those calls, but when it
is run over continuous field recordings it also fires on the many non-target
sounds that dominate a rainforest soundscape --- insects, birds, wind, rain,
and the calls of non-target animals --- none of which are well represented in
the original training set. The result is a detector with a high false-positive
rate, and verifying every detection by hand reintroduces exactly the
manual-listening bottleneck that automation was meant to remove.

This method addresses that problem for two primate species recorded at
Makokou, Gabon: the putty-nosed monkey (\textit{Cercopithecus nictitans}),
whose repertoire includes acoustically distinct ``hack'', ``kek'', and
``pyow'' call types \cite{arnold2006}, and \textit{Colobus guereza}. Rather
than only training a classifier, we wrap it in (i) a detection stage with
domain-appropriate post-processing, (ii) an automatic three-filter cleanup
that flags likely false positives using signals independent of the detector
itself, and (iii) an iterative hard-negative mining loop that recycles
confirmed false positives into the background class and retrains. The
motivation is to obtain trustworthy detections from long recordings while
keeping the amount of manual listening close to zero, and to do so with a
workflow that other researchers can re-run on their own species by supplying
new reference clips and editing a single configuration file.

%% =====================================================================
%% METHOD DETAILS -- enough detail to reproduce the method. No word limit.
%% =====================================================================
\section*{Method details}

The pipeline has five stages --- preprocessing, training, detection, automatic
cleanup, and iterative retraining. All paths, species definitions, and
hyperparameters live in a single \texttt{config.py} so the workflow can be
re-targeted to new species or sites without touching the code.

\subsection*{1. Study area, recordings, and reference clips}

Field recordings are continuous WAV files collected at fixed acoustic stations
(denoted IPA1ST, IPA2ST, \dots; one subfolder per recording date) in
rainforest near Makokou, Gabon. Because the target species are diurnal and
most vocally active in the morning, only recordings whose start time falls
within a configurable window (default 05:30--10:30 local time, parsed from the
timestamp in each filename) are processed; this removes large stretches of
audio in which the species are unlikely to call and reduces both compute and
false positives.

The model is trained on short reference clips organized into three classes:
\begin{itemize}
\item \textbf{Cernic} --- all putty-nosed monkey call types (hack, kek, pyow)
pooled into a single class. Four source folders (one per call type, plus
field-confirmed recoveries) are mapped to the single ``Cernic'' label via
\texttt{config.SPECIES\_FOLDERS}, so that merging happens at training time and
the model directly outputs a unified putty-nosed decision;
\item \textbf{Colobus\_guereza} --- \textit{Colobus guereza} calls;
\item \textbf{Background} --- a pooled negative class assembled from ambient
forest noise, calls of non-target animals present at the site (e.g.\
\textit{Cercocebus torquatus}, \textit{Pan troglodytes}), and previously
misclassified clips recycled through the hard-negative mining loop.
\end{itemize}

\noindent All training clips underwent human verification. For Colobus, 69
mislabelled clips (containing only bird calls due to mechanical slicing of
longer recordings) were identified and removed. For the Background class, a
label-quality audit of the entire auto-flagged false-positive pool (4669
files) was conducted using model-confidence-based rescoring and manual review;
six clips confirmed as real Cernic calls were recovered and moved to a
dedicated \texttt{field\_confirmed} training folder.

\subsection*{2. Audio preprocessing and mel-spectrogram representation}

All audio is resampled to 44.1\,kHz and standardized to a fixed 2.0\,s
duration. Clips shorter than the target (e.g.\ hack, kek, and pyow calls,
typically $\sim$0.3\,s) are embedded at a random position within a 2\,s
segment of real background noise randomly selected from the Background pool;
this eliminates the silence shortcut that zero-padding would introduce and
forces the model to distinguish the call from genuine ambient sound. Clips
longer than 2\,s (e.g.\ putty-nosed 5\,s field recordings) are trimmed to the
deterministic 2\,s window with the highest RMS energy
(\texttt{find\_loudest\_window}), ensuring that the segment reviewed by a
human annotator is identical to the segment used for training. Each 2\,s
waveform is converted to a mel-spectrogram with librosa \cite{mcfee2015} using
a 2048-sample FFT, 512-sample hop, 128 mel bands, and a 20--8000\,Hz analysis
range, and is expressed in decibels relative to its per-clip maximum. The
spectrogram is min--max normalized to [0, 255], resized to
224\,$\times$\,224 with bilinear interpolation, replicated across three
channels to form an RGB image, and finally scaled to [0, 1] for input to the
network. Using the same conversion at training and detection time keeps the
two feature distributions comparable.

\subsection*{3. Data augmentation}

To expand the reference set and make the classifier robust to field
conditions, each target-class spectrogram is expanded into seven training
examples (a 7$\times$ multiplier): one unaltered copy; three copies mixed with
randomly chosen background spectrograms at a signal-to-noise ratio drawn
uniformly from $-5$ to 10\,dB; one time-axis crop and one frequency-axis crop
(each removing 10--30\,\% from a randomly chosen edge, then resized back to the
original shape); and one frequency translation of $\pm$20 mel bins.
Background-class clips are not augmented, so that the negative class reflects
the true diversity of ambient sound rather than synthetic variants. The
augmentation operations are spectrogram-domain analogues of standard
SpecAugment-style masking and mixing \cite{park2019}.

\subsection*{4. Model architecture}

The classifier couples a VGG19 convolutional backbone \cite{simonyan2015}
pretrained on ImageNet \cite{deng2009} with a custom temporal-frequency CRNN
head designed to exploit both the spectral location and the temporal rhythm of
primate calls. The motivation is that Colobus roars concentrate energy around
$\sim$512\,Hz, putty-nosed calls around $\sim$1\,kHz, and common confusers
(birds, insects) above $\sim$2\,kHz; encoding \emph{where} in frequency the
energy sits and \emph{when} it occurs lets the model discriminate species that
a global pooling layer would conflate.

The feature map from \texttt{block4\_conv4} of VGG19 (spatial size
28\,$\times$\,28, 512 channels) is extracted before the final VGG19 block. The
frequency axis (height) is split into four equal bands (each 7 bins), so that
each band covers a contiguous portion of the mel-frequency range. Within each
band, the frequency dimension is collapsed (average pooling over the 7 bins),
yielding a time-$\times$-channel sequence of shape 28\,$\times$\,512. Each
band is then processed by its own one-dimensional convolutional stream ---
Conv1D(128 filters, kernel size 3, same padding) followed by batch
normalization and ReLU --- producing four sequences of shape
28\,$\times$\,128. These are concatenated along the channel axis
(28\,$\times$\,512) and passed through a cross-band Conv1D(256, kernel 3,
same padding) + batch normalization + ReLU layer that allows the network to
learn interactions between frequency bands.

The resulting 28\,$\times$\,256 sequence is fed into a bidirectional LSTM with
128 units per direction (return sequences enabled, recurrent dropout 0.3),
which captures temporal dependencies such as the rhythmic hack--hack or
pyow--hack patterns characteristic of putty-nosed calling bouts. The LSTM
output sequence is summarized by concatenating global max pooling and global
average pooling over the time axis, producing a 512-dimensional vector that
retains both peak activation and average energy. This vector passes through
two dense layers (256 units, ReLU, dropout 0.5; then 128 units, ReLU, dropout
0.5) and a final three-way softmax output. The full model has approximately
12.6 million trainable parameters.

\subsection*{5. Model training}

Training proceeds in two stages, following a standard transfer-learning
protocol of head training followed by partial fine-tuning.

\paragraph{Stage 1 --- Head training.}
The entire VGG19 convolutional base is frozen and only the temporal-frequency
CRNN head is trained. The optimizer is Adam with an initial learning rate of
$1 \times 10^{-4}$, and the loss is sparse categorical cross-entropy.

\paragraph{Stage 2 --- Partial fine-tuning.}
The last two convolutional blocks of VGG19 (\texttt{block4} and
\texttt{block5}) are unfrozen while earlier blocks remain frozen to preserve
low-level features. The learning rate is reduced to $1 \times 10^{-5}$ to
avoid catastrophic forgetting.

\paragraph{Shared training settings.}
Both stages use the augmented dataset split 80/20 into training and validation
sets with stratification by class and a fixed random seed for
reproducibility. Balanced class weights (inversely proportional to class
frequency) compensate for the larger Background class. Training runs for up to
50 epochs per stage with a batch size of 32, governed by three callbacks:
early stopping on validation accuracy (patience 10, best weights restored),
model checkpointing on best validation accuracy, and learning-rate reduction
on plateau (factor 0.5, patience 5, floor $1 \times 10^{-7}$). The best
checkpoint is saved as \texttt{best\_model.h5}. Training concludes with a
per-class accuracy report and a confusion matrix on the validation set.

\subsection*{6. Sliding-window detection}

A trained model is applied to each long recording with a 2.0\,s window
advanced in 1.0\,s steps (50\,\% overlap); a final window is anchored at the
end of the file so trailing audio is not dropped. Every window is converted to
a mel-spectrogram image and classified. Because the three putty-nosed call
types are already merged into the single Cernic class at training time, no
softmax summation across call types is needed at detection time ---
\texttt{DETECTION\_GROUPS} is simply an identity mapping (Cernic $\to$ Cernic,
Colobus $\to$ Colobus, Background $\to$ Background). The window is assigned to
its highest-scoring class and is kept as a detection if that class is not
Background and its confidence meets a threshold (default 0.4). Overlapping
detections of the same class are then reduced by per-class non-maximum
suppression at an intersection-over-union threshold of 0.5. Per-file
detections (start time, end time, species, confidence) are written to CSV.

\subsection*{7. Automatic false-positive cleanup}

The detector's raw output contains many false positives from sounds absent
from the training set. Three independent filters, each exploiting a different
signal, are run over the saved detections; their agreement determines how a
detection is treated. For each detection a 2\,s clip is extracted from the
source recording.

\paragraph{(a) Mahalanobis out-of-distribution filter}
Using a penultimate dense layer as a feature extractor, per-class feature means
and a pooled, ridge-regularized inverse covariance are estimated from the
original (non-augmented) training clips \cite{lee2018}. For each detection the
Mahalanobis distance to its candidate class cluster is computed. A real call
lies close to its training cluster, whereas an unfamiliar sound does not.
Because field clips are systematically noisier than clean training clips,
thresholds calibrated on training data flag almost everything; we therefore
calibrate the cutoff on the \textbf{detection distances themselves}, flagging,
per species, the detections beyond the 95th percentile of the observed
distance distribution (a configurable percentile). This makes the filter a
relative-outlier detector that is robust to the train-to-field domain shift.

\paragraph{(b) YAMNet cross-check}
Each clip is resampled to 16\,kHz and passed to YAMNet, a network trained on
Google's AudioSet ontology of 521 everyday sound classes
\cite{gemmeke2017,hershey2017}. A detection is flagged when YAMNet's top class
is one of a curated set of non-primate sounds (birds, insects, frogs, wind,
rain, water, mechanical and human sounds, etc.); generic classes consistent
with a primate call (e.g.\ ``Animal'', ``Wild animals'') are deliberately left
unflagged. Because AudioSet has no monkey-specific class, this filter is used
only to \emph{reject} clips that are confidently something else, not to
confirm primates.

\paragraph{(c) Temporal-isolation filter}
Primates tend to call in bouts, so a genuine call usually has same-species
neighbours nearby in time. A detection with no other detection of the same
species within $\pm$30\,s (configurable) is flagged as temporally isolated.

A detection is considered \textbf{clean} only when no filter flags it,
\textbf{suspicious} when at least one filter flags it, and a \textbf{strong
false positive} when at least two filters agree. Clean and suspicious
detections are written to separate CSVs (the suspicious file carries a
human-readable \texttt{flag\_reason} column), and strong false positives are
exported as labelled WAV clips for the next stage.

\subsection*{8. Iterative hard-negative mining}

The exported strong false positives are the sounds the detector most needs to
learn to ignore. They are added to the Background class and the model is
retrained from the preprocessing stage, so that on the next pass the network
is far less likely to fire on those sounds. Repeating the detect $\rightarrow$
clean $\rightarrow$ fold-in $\rightarrow$ retrain loop a few times (typically
three to five iterations) progressively drives down the false-positive rate
without ever requiring exhaustive manual annotation of the field recordings.

%% =====================================================================
%% METHOD VALIDATION -- data that validate the method (no Results/Discussion).
%% =====================================================================
\section*{Method validation}

We report two levels of validation: (i) performance on a held-out validation
split drawn from the same curated clip pool as training, and (ii) a field test
on continuous recordings from one acoustic station not used during development.

\paragraph{Validation-set performance (V8 model).}
The validation set contains 2713 clips stratified across the three classes.
Table~\ref{tab:val} summarizes per-class results.

\begin{table}[h]
\centering
\caption{Per-class validation-set performance of the final (V8) model.}
\label{tab:val}
\begin{tabular}{lrrr}
\toprule
Class & Samples & Accuracy (\%) & Mean confidence \\
\midrule
Cernic           & 492  & 97.15 & 0.969 \\
Colobus\_guereza & 864  & 99.88 & 0.998 \\
Background       & 1357 & 99.34 & 0.992 \\
\midrule
\textbf{Overall} & \textbf{2713} & \textbf{99.12} & --- \\
\bottomrule
\end{tabular}
\end{table}

\noindent Overall accuracy is 99.12\% with a loss of 0.036. Notably, the
confusion matrix shows \textbf{zero cross-species confusion} between Cernic
and Colobus. All errors involve the Background class: 14 Cernic clips are
conservatively classified as Background (mostly weak or brief hacks and pyows
embedded in noise), while 6 Background clips are misclassified as Cernic and 3
as Colobus.

\paragraph{Field test --- IPA19ST.}
The V8 model was applied to 40 continuous recordings from station IPA19ST.
Of 19 total detections, 17 were attributed to Cernic and 2 to Colobus:
\begin{itemize}
\item \textbf{Cernic}: 14 true positives out of 17 detections
(precision 82.4\%).
\item \textbf{Colobus}: 0 true positives out of 2 detections (precision 0\%;
both were overlapping events from the sliding window).
\end{itemize}

\paragraph{Model-version progression on IPA19ST.}
The improvement across architecture iterations is striking. The V6 model
(frequency-band pooling replacing global average pooling) produced 39
detections, all false positives (precision 0\%). The V7 model (temporal CRNN
head with Conv1D + BiLSTM but no frequency-band split) yielded 3 detections
with 1 true positive. The final V8 model (temporal-frequency CRNN) recovered
Cernic recall (14 true positives) while keeping false positives manageable (5
total FP out of 19 detections).

\paragraph{Pending validation.}
Field-test results for station IPA20ST are pending and will be reported in a
future update.

%% =====================================================================
\section*{Limitations}

\begin{itemize}
\item The training set contains no real Colobus field recordings from the
Makokou site; all Colobus reference clips are sourced from curated databases.
This domain mismatch likely explains the poor Colobus precision on field data
(0 of 2 detections on IPA19ST). Collecting site-specific Colobus recordings
is a priority for future work.
\item The auto-cleanup pipeline introduces a circular bias: hard negatives
are scored by the same model that produced them, so genuine calls with low
model confidence may be systematically hidden. This risk is mitigated by a
mining-confidence audit in which high-confidence clips in the false-positive
pool are manually reviewed, and confirmed calls are recovered into the
training set.
\item The model is conservative on weak or short calls. On the validation
set, 14 Cernic clips (2.8\%) --- mostly brief hacks and pyows embedded in
loud ambient noise --- are misclassified as Background. This trade-off
favours precision over recall.
\item Re-targeting to new species requires curated reference clips and edits
to the configuration; performance on rare or quiet calls with few reference
examples is not guaranteed.
\item The YAMNet cross-check has no primate-specific class and only rejects
clips it can confidently attribute to a non-primate sound; tropical-forest
sounds outside AudioSet's training distribution may be mislabelled, so the
filter is conservative by design.
\item The Mahalanobis filter assumes that the majority of detections for a
species are genuine, so that outliers are informative; if a species'
detections are overwhelmingly false positives, the relative-outlier
calibration is less effective.
\item The temporal-isolation filter assumes bout-calling; genuinely isolated
or rare calls can be flagged, and very dense calling can mask true isolation.
\item A fixed 2\,s analysis window may truncate or split long vocalizations
(e.g.\ \textit{Colobus} roars), and the time-of-day filter (05:30--10:30) is
site- and species-specific and must be adjusted elsewhere.
\item The method yields detections, not calibrated abundance or density
estimates; converting counts to ecological metrics requires additional
assumptions.
\end{itemize}

%% =====================================================================
%% ETHICS STATEMENTS -- keep only the statements relevant to your work.
%% =====================================================================
\section*{Ethics statements}

This study involved only the passive recording and analysis of environmental
audio from free-ranging animals in their natural habitat. It included no
animal experiments and no capture, handling, or manipulation of animals; no
human subjects; and no data collected from social-media platforms. [If
applicable, add: field recordings were collected under permit [number] issued
by [authority], in accordance with the relevant institutional and national
guidelines for non-invasive field research.]

%% =====================================================================
%% CREDIT AUTHOR STATEMENT
%% =====================================================================
\section*{CRediT author statement}

\textbf{Moshi Fu:} Conceptualization, Methodology, Software, Validation,
Formal analysis, Data curation, Writing -- original draft, Visualization.
\textbf{[Co-author Name]:} [Supervision, Resources, Writing -- review \&
editing, Funding acquisition].

%% =====================================================================
\section*{Acknowledgments}

[Acknowledge field teams, station maintainers, and anyone who contributed
reference annotations but does not meet authorship criteria.] This research did
not receive any specific grant from funding agencies in the public,
commercial, or not-for-profit sectors.

%% =====================================================================
%% DECLARATION OF INTERESTS -- keep BOTH statements; tick the correct box.
%% =====================================================================
\section*{Declaration of interests}

\noindent$\boxtimes$~The authors declare that they have no known competing
financial interests or personal relationships that could have appeared to
influence the work reported in this paper.

\medskip
\noindent$\square$~The authors declare the following financial
interests/personal relationships which may be considered as potential
competing interests:

%% =====================================================================
\section*{Supplementary material and/or additional information [OPTIONAL]}

[Optional. List any supplementary files here, e.g.\ example reference clips,
the \texttt{config.py} used, or sample detection CSVs. Delete this section if
not used.]

%% =====================================================================
\begin{thebibliography}{99}

\bibitem{stowell2022}
D. Stowell, Computational bioacoustics with deep learning: a review and
roadmap, PeerJ 10 (2022) e13152. \url{https://doi.org/10.7717/peerj.13152}

\bibitem{simonyan2015}
K. Simonyan, A. Zisserman, Very deep convolutional networks for large-scale
image recognition, in: International Conference on Learning Representations
(ICLR), 2015. \url{https://doi.org/10.48550/arXiv.1409.1556}

\bibitem{deng2009}
J. Deng, W. Dong, R. Socher, L.-J. Li, K. Li, L. Fei-Fei, ImageNet: a
large-scale hierarchical image database, in: IEEE Conference on Computer
Vision and Pattern Recognition (CVPR), 2009, pp. 248--255.
\url{https://doi.org/10.1109/CVPR.2009.5206848}

\bibitem{gemmeke2017}
J.F. Gemmeke, D.P.W. Ellis, D. Freedman, A. Jansen, W. Lawrence, R.C. Moore,
M. Plakal, M. Ritter, Audio Set: an ontology and human-labeled dataset for
audio events, in: IEEE International Conference on Acoustics, Speech and Signal
Processing (ICASSP), 2017, pp. 776--780.
\url{https://doi.org/10.1109/ICASSP.2017.7952261}

\bibitem{hershey2017}
S. Hershey, S. Chaudhuri, D.P.W. Ellis, J.F. Gemmeke, A. Jansen, R.C. Moore,
M. Plakal, D. Platt, R.A. Saurous, B. Seybold, M. Slaney, R.J. Weiss,
K. Wilson, CNN architectures for large-scale audio classification, in: IEEE
International Conference on Acoustics, Speech and Signal Processing (ICASSP),
2017, pp. 131--135. \url{https://doi.org/10.1109/ICASSP.2017.7952132}

\bibitem{lee2018}
K. Lee, K. Lee, H. Lee, J. Shin, A simple unified framework for detecting
out-of-distribution samples and adversarial attacks, in: Advances in Neural
Information Processing Systems (NeurIPS) 31, 2018, pp. 7167--7177.
\url{https://doi.org/10.48550/arXiv.1807.03888}

\bibitem{mcfee2015}
B. McFee, C. Raffel, D. Liang, D.P.W. Ellis, M. McVicar, E. Battenberg,
O. Nieto, librosa: audio and music signal analysis in Python, in: Proceedings
of the 14th Python in Science Conference (SciPy), 2015, pp. 18--25.
\url{https://doi.org/10.25080/Majora-7b98e3ed-003}

\bibitem{park2019}
D.S. Park, W. Chan, Y. Zhang, C.-C. Chiu, B. Zoph, E.D. Cubuk, Q.V. Le,
SpecAugment: a simple data augmentation method for automatic speech
recognition, in: Interspeech, 2019, pp. 2613--2617.
\url{https://doi.org/10.21437/Interspeech.2019-2680}

\bibitem{kahl2021}
S. Kahl, C.M. Wood, M. Eibl, H. Klinck, BirdNET: a deep learning solution for
avian diversity monitoring, Ecol. Inform. 61 (2021) 101236.
\url{https://doi.org/10.1016/j.ecoinf.2021.101236}

\bibitem{dufourq2022}
E. Dufourq, C. Batist, R. Foquet, I. Durbach, Passive acoustic monitoring of
animal populations with transfer learning, Ecol. Inform. 70 (2022) 101688.
\url{https://doi.org/10.1016/j.ecoinf.2022.101688}

\bibitem{arnold2006}
K. Arnold, K. Zuberb\"uhler, Semantic combinations in primate calls, Nature
441 (2006) 303. \url{https://doi.org/10.1038/441303a}

\end{thebibliography}

\end{document}
